\documentclass[10pt,a4paper]{amsart}
\usepackage[margin=19mm]{geometry}
\usepackage{amsmath,amssymb,mathtools}
\usepackage[scr=rsfs,cal=euler]{mathalfa}
\usepackage{xfrac}
\usepackage[unicode,hidelinks,bookmarks=false]{hyperref}
\newcommand{\ie}{i.e.\ }
\newcommand{\cf}{cf.\ }
\newcommand{\resp}{respectively\ }
\newcommand{\Subsection}{\subsection}
\providecommand{\nobreakdash}{\nobreak\hbox{-}}
\setlength{\emergencystretch}{2em}
\multlinegap0pt
\usepackage{bm}
\usepackage{bbm}
\usepackage{dsfont}
\usepackage{xspace}
\usepackage{cancel}
\usepackage{mathtools}
\usepackage{amsthm}
\makeatletter
\@ifundefined{theorem}{\newtheorem{theorem}{Theorem}}{}
\makeatother
\newcommand{\Aa}{{\mathcal A}}
\newcommand{\Ee}{{\mathcal E}}
\newcommand{\Ss}{{\mathcal S}}
\newcommand{\Z}{\mathbb Z}
\title[Contact flexibility and rigidity for toric Gorenstein prequa]{Contact flexibility and rigidity for toric Gorenstein prequantizations and Ehrhart theory of toric diagrams}
\author{Miguel Abreu, Leonardo Macarini, Ant\'onio Rocha-Neves}
\date{Source: arXiv:2604.23848v1 (CC BY 4.0)}
\begin{document}
\maketitle

\noindent\emph{Source and licence.} Contact flexibility and rigidity for toric Gorenstein prequantizations and Ehrhart theory of toric diagrams. Version arXiv:2604.23848v1 (CC BY 4.0), distributed under the Creative Commons Attribution 4.0 International licence, as recorded in the arXiv licence metadata for this exact version and shown on the abstract page https://arxiv.org/abs/2604.23848v1. Full rights notice: licenses/arxiv-2604.23848-CC-BY-4.0.txt.\par\smallskip

\noindent\textit{Editorial scope.} Counted unit: the Theorem whose source label is \texttt{thm:bipysimplex2} in the article (source0.tex lines 2476--2480, source label \texttt{thm:bipysimplex2}), delivered together with the article's own complete original proof of it (source0.tex lines 2481--2536, one \texttt{proof} environment of 56 lines). No mathematics is written, repaired or re-derived: statement and proof are the article's own text line for line. Required context retained uncounted: thm:uniqueness (lines 2455--2458). Every definition, display and label of those lines is retained unchanged. Omitted from this package and not redistributed: 1--2454, 2459--2475, 2537--2598. Nothing inside a retained excerpt is omitted or altered: every retained body is delivered exactly as the article's lines are written, so no omission receipt of any kind accompanies this package. Citations: the retained excerpts carry no citation command at all, so no bibliography entry of the article is transcribed and no reference is redistributed. Reference adaptation: none was needed and none was made; every reference inside the delivered unit resolves inside it, so no unresolved reference or citation is produced. Printed slips kept literal: no printed word, symbol, hypothesis or number of the counted statement or of its proof was altered. Layout and class: the article's own class is \texttt{amsart} with packages including a4wide, inputenc, fontenc, amsmath, amssymb, amsthm, fancyhdr, graphicx, enumerate, cite, appendix, hyperref, xy, latexsym; the wrapper is the lane's pinned \texttt{amsart} editorial wrapper at 10pt on A4, which restates the theorem-like environments the retained text needs and adds the article's own macro definitions that the retained text needs (\texttt{\textbackslash Aa}, \texttt{\textbackslash Ee}, \texttt{\textbackslash Ss}, \texttt{\textbackslash Z}), with extra wrapper packages (bm, bbm, dsfont, xspace, cancel, mathtools). The delivered numbering is the unit's own. Assets: no retained span contains a figure environment, a graphics-inclusion command, a PDF-page inclusion, a table or a TikZ picture; no image file is copied into this package, the package declares no asset dependency and no image file is redistributed with it. Licence: version arXiv:2604.23848v1 of this article is distributed under the Creative Commons Attribution 4.0 International licence, as recorded in the arXiv licence metadata for this exact version and shown on the abstract page https://arxiv.org/abs/2604.23848v1. A full rights notice is delivered at licenses/arxiv-2604.23848-CC-BY-4.0.txt. Characters: every retained excerpt is pure ASCII, so the delivered engine, XeLaTeX with its default Latin Modern text font, reports no missing character.
\begin{theorem} \label{thm:uniqueness}
For integers $k_1, k_2$ such that $0 \le k_1, k_2 < n$ and $k_1 \equiv k_2 \equiv n \pmod 2$, the prequantized 
toric diagrams $\mathcal{D}_{k_1}$ and $\mathcal{D}_{k_2}$ are unimodularly equivalent if and only if $k_1 = k_2$.
\end{theorem}

\begin{theorem} \label{thm:bipysimplex2}
Let $\Ss$ be a $n$-dimensional toric diagram. If its $h^*$-polynomial is $h_{\Ss}^\ast (z) = 1 + z + z^2 + \dots + z^{n-1}$, 
then $\Ss$ is unimodularly equivalent to $\mathcal{D}_k$ for a unique $k\in\Z$ satisfying $0 \le k < n$ and 
$k \equiv n \pmod 2$.
\end{theorem}

\begin{proof}
From the given $h^\ast$-polynomial, the normalized volume of $\Ss$ is $h_{\Ss}^\ast(1) = n$. 
The total number of lattice points is $L_{\Ss}(1) = h_1^\ast + n + 1 = n+2$. Since the degree of 
$h_{\Ss}^\ast$ is $n-1$ and it is palindromic, $\Ss$ is Gorenstein of index $2$ thus contains no interior 
lattice points. Because $\Ss$ is a toric diagram, its facets are primitive, so all $n+2$ lattice points are vertices.
		
Next, we determine the number of facets. The standard expansion of the Ehrhart polynomial dictates that the 
coefficient of $t^{n-1}$ in $L_{\Ss}(t)$ is
\[
\frac{1}{2} \sum_{F \in \mathcal{F}(\Ss)} \text{Vol}_{n-1}(F). 
\]
Because every facet $F$ of a toric diagram is a unimodular simplex, the relative volume 
$\text{Vol}_{n-1}(F) = \frac{1}{(n-1)!}$, making this coefficient exactly $\frac{f_{n-1}}{2(n-1)!}$. Using our specific 
$h^\ast$-polynomial, the Ehrhart polynomial is $L_{\Ss}(t) = \sum_{j=0}^{n-1} \binom{t+n-j}{n}$, where the coefficient 
of $t^{n-1}$ is $\frac{n}{(n-1)!}$. Equating these gives a facet count of $f_{n-1} = 2n$.
		
Any $n$-dimensional polytope with $n+2$ vertices admits a unique (up to scaling) affine dependency 
$\sum \lambda_i \bm{v}_i = 0$ with $\sum \lambda_i = 0$. Because every facet of $\Ss$ is a simplex, no subset of 
$n+1$ vertices can be affinely dependent. This strictly guarantees that $\lambda_i \neq 0$ for all $i$, meaning every vertex 
participates in the dependency. This dependency therefore partitions the vertices into exactly two disjoint sets 
$I = \{\bm{v}_i : \lambda_i > 0\}$ and $J = \{\bm{v}_j : \lambda_j < 0\}$. A subset of vertices forms a facet if and only if it is 
obtained by removing exactly one vertex from $I$ and one from $J$. Thus, the number of facets is given by the product 
$|I| \cdot |J|$. We are led to the system
\[ 
|I| + |J| = n+2 \quad \text{and} \quad |I| \cdot |J| = 2n\,,
\]
whose unique integer solution is $\{|I|, |J|\} = \{n, 2\}$. A polytope whose vertices are partitioned in this way is, 
by definition, a bipyramid over a $(n-1)$-simplex. Let $\Ee = \{\bm{v}_1, \dots, \bm{v}_n\}$ be the base vertices and 
$A = \{\bm{y}_1, \bm{y}_2\}$ be the apices.
		
The bipyramid $\Ss$ is triangulated by $n$ lattice simplices $S_i = \text{conv}((\Ee \setminus \{\bm{v}_i\}) \cup \Aa)$,
therefore the sum of their normalized volumes must be $n$. Since each $S_i$ is a lattice simplex, its normalized volume must 
be an integer greater than or equal to $1$. This forces the normalized volume of $S_i$ to be $1$ for all $1 \le i \le n$. The equality 
of these volumes implies that the coefficients of the base vertices in the affine dependency are identical. The dependency relation 
is thus
\[
\sum_{i=1}^n \bm{v}_i = \mu_1 \bm{y}_1 + \mu_2 \bm{y}_2 \,,
\]
where $\mu_1 + \mu_2 = n$. By setting $k = \mu_2 - \mu_1$, we find $\mu_1 = \frac{n-k}{2}$ and $\mu_2 = \frac{n+k}{2}$. 
Up to swapping the apices $\bm{y}_1$ and $\bm{y}_2$, we may assume $0 \le k < n$.
		
Finally, because $\text{Vol}(S_n) = 1$, the vectors $\{\bm{v}_1 - \bm{y}_1, \dots, \bm{v}_{n-1} - \bm{y}_1, \bm{y}_2 - \bm{y}_1\}$ 
form a $\mathbb{Z}$-basis for the lattice. Let $\phi$ be the affine unimodular transformation that maps this basis to the basis 
formed by the corresponding vertices of $\mathcal{D}_k$:
\[ 
\phi(\bm{y}_1) = \bm{a}_n, \quad \phi(\bm{y}_2) = \bm{a}_{n+1}, \quad \phi(\bm{v}_i) = \bm{a}_i \text{ for } 1 \le i \le n-1\,.
\]
Applying $\phi$ to our dependency relation isolates $\phi(\bm{v}_n)$
\[
\sum_{i=1}^{n-1} \bm{a}_i + \phi(\bm{v}_n) = \frac{n-k}{2} \bm{a}_n + \frac{n+k}{2} \bm{a}_{n+1} \implies 
\phi(\bm{v}_n) = \left(-1, \dots, -1, -k, \frac{n+k}{2}\right) \,.
\]
For $\phi(\bm{v}_n)$ to be an integral lattice point, the coordinate $\frac{n+k}{2}$ must be an integer, which forces the condition 
$k \equiv n \pmod 2$. Therefore, $\Ss$ is unimodularly equivalent to $\mathcal{D}_k$. Uniqueness of $k$ is a consequence
of Theorem~\ref{thm:uniqueness}.
\end{proof}

\end{document}
